% ============================================================
%  Virtual Ambulance Emergency Response Simulation
%  Final Course Project Report — SBEG353 / SBES140
%  Spring 2026 | Cairo University
%  Two-column IEEE-style layout
% ============================================================

\documentclass[10pt, twocolumn]{article}

% ── Packages ────────────────────────────────────────────────
\usepackage[a4paper, top=2cm, bottom=2cm, left=1.5cm, right=1.5cm,
            columnsep=0.6cm]{geometry}
\usepackage{times}          % Times New Roman font
\usepackage{graphicx}       % \includegraphics
\usepackage{float}          % [H] figure placement
\usepackage{caption}        % caption formatting
\usepackage{subcaption}     % subfigures
\usepackage{booktabs}       % professional tables
\usepackage{amsmath}        % math
\usepackage{hyperref}       % clickable links
\usepackage{enumitem}       % list control
\usepackage{titlesec}       % section title formatting
\usepackage{fancyhdr}       % header / footer
\usepackage{multicol}       % span columns when needed
\usepackage{array}
\usepackage{xcolor}
\usepackage{parskip}
\usepackage{microtype}

% ── Caption style ───────────────────────────────────────────
\captionsetup{font=small, labelfont=bf, skip=4pt}

% ── Section formatting ──────────────────────────────────────
\titleformat{\section}{\normalfont\large\bfseries}{\thesection.}{0.5em}{}
\titleformat{\subsection}{\normalfont\normalsize\bfseries}{\thesubsection.}{0.5em}{}
\titlespacing*{\section}{0pt}{8pt}{4pt}
\titlespacing*{\subsection}{0pt}{6pt}{3pt}

% ── Header / Footer ─────────────────────────────────────────
\pagestyle{fancy}
\fancyhf{}
\lhead{\small SBEG353/SBES140 — Computer Graphics and Visualization, Spring 2026}
\rhead{\small Cairo University}
\cfoot{\thepage}
\renewcommand{\headrulewidth}{0.4pt}

% ── Hyperref ────────────────────────────────────────────────
\hypersetup{colorlinks=true, linkcolor=blue, urlcolor=blue, citecolor=blue}

% ============================================================
\begin{document}

% ── Title block (full width) ────────────────────────────────
\twocolumn[
  \begin{@twocolumnfalse}
    \begin{center}
      {\Large \textbf{Virtual Ambulance Emergency Response Simulation in Blender}} \\[6pt]
      {\normalsize Final Course Project — SBEG353 / SBES140 | Spring 2025/2026} \\[4pt]
      {\small
        \textbf{Author 1}\textsuperscript{1},
        \textbf{Author 2}\textsuperscript{1},
        \textbf{Author 3}\textsuperscript{1},
        \textbf{Author 4}\textsuperscript{1},
        \textbf{Author 5}\textsuperscript{1} \\[2pt]
        \textsuperscript{1}Systems and Biomedical Engineering Department,
        Cairo University, Giza, Egypt \\[2pt]
        \texttt{\{student\_id1, student\_id2, student\_id3,
                  student\_id4, student\_id5\}@student.cu.edu.eg}
      }
    \end{center}
    \vspace{6pt}
    \hrule
    \vspace{6pt}
    % ── Abstract ──────────────────────────────────────────
    \begin{center}
      \textbf{Abstract}
    \end{center}
    \noindent
    This report presents the design and implementation of a Virtual Ambulance
    Emergency Response Simulation developed in Blender as the final project for
    the Computer Graphics and Visualization course (SBEG353/SBES140).
    The simulation models a realistic urban emergency medical response workflow,
    encompassing a multi-scene virtual environment that includes city streets,
    hospital interiors, and accident sites.
    The project integrates core computer-graphics concepts, including
    environment modelling, UV mapping, texture application, skeletal character
    animation, physics-based rendering, advanced lighting, and collision handling.
    Three distinct emergency-response scenarios are implemented: a traffic
    collision, a hospital reception handover, and a patient evacuation by
    stretcher.
    A fire-emergency overlay and ambulance dispatch sequence further demonstrate
    dynamic lighting and particle-system capabilities.
    The report documents the modelling workflow, animation pipeline, lighting
    setup, and physics simulations employed, as well as the contribution of each
    team member.
    \\[6pt]
    \textbf{Keywords:} Blender, computer graphics, skeletal animation,
    emergency simulation, physics-based rendering, collision detection,
    ambulance, virtual environment.
    \vspace{6pt}
    \hrule
    \vspace{10pt}
  \end{@twocolumnfalse}
]

% ============================================================
\section{Introduction}
% ============================================================

Virtual simulation has become a powerful tool in medical education and emergency
preparedness, enabling practitioners to rehearse complex scenarios without risk to
real patients~\cite{caballero2018, lu2021}.
This project applies computer-graphics principles to construct a Virtual Ambulance
Emergency Response Simulation that immerses the viewer in a sequence of
life-saving events, from the initial traffic accident through on-scene first aid,
ambulance transport, and hospital reception.

The simulation was entirely built in Blender~3.x and targets the following learning
objectives defined in the project specification:
(i)~virtual environment creation,
(ii)~hero object modelling,
(iii)~skeletal character animation,
(iv)~emergency response scenarios,
(v)~lighting and rendering,
(vi)~collision and object interaction, and
(vii)~physics-based rendering.

The remainder of this report is structured as follows.
Section~\ref{sec:env} describes the virtual environment.
Section~\ref{sec:hero} covers hero-object design.
Section~\ref{sec:char} details character rigging and animation.
Section~\ref{sec:scenarios} presents the three emergency scenarios.
Section~\ref{sec:lighting} discusses lighting and rendering.
Section~\ref{sec:collision} addresses collision handling.
Section~\ref{sec:physics} explains the physics-based simulation.
Section~\ref{sec:contributions} lists member contributions.
Section~\ref{sec:conclusion} concludes the report.

% ============================================================
\section{Virtual Environment Creation}
\label{sec:env}
% ============================================================

\subsection{Urban Street and City Block}

The primary outdoor environment represents a dense urban district modelled from
reference photographs of mid-rise commercial blocks.
High-rise glass office towers with grid-pattern façades line both sides of a
multi-lane intersection equipped with pedestrian crosswalks, street lamps, traffic
signals, and planted trees.
All building surfaces utilise PBR materials: glass panels employ a specular BSDF
combined with a glossy node to achieve realistic reflections, while concrete pavements
are textured with a diffuse-rough BSDF to simulate worn stone.

% ── Figure: City overview ──────────────────────────────────
\begin{figure}[H]
  \centering
  \includegraphics[width=\linewidth]{IMG-20260618-WA0070.jpg}
  % Replace with your actual filename if different
  \caption{Aerial overview of the modelled urban city block showing high-rise
           buildings, intersections, street lamps, and planted trees.}
  \label{fig:city_overview}
\end{figure}

\subsection{Hospital Exterior and Interior}

A dedicated hospital building was constructed with a prominent red cross emblem
on its façade and ground-floor glazed entrance.
Inside, the reception area features a curved wooden desk labelled
\textit{``Welcome / Reception''}, office chairs, flat-panel monitors, a
\textsc{Caution: Wet Floor} sign, a wheelchair, and department signage
(\textit{Intensive Care Unit}, \textit{Emergency Department}, \textit{Radiology}).
Material variety was deliberately maximised: the reception counter uses a wood
shader, the floor tiles a specular terrazzo material, and the glass partition a
transparent BSDF.

% ── Figure: Hospital exterior ──────────────────────────────
\begin{figure}[H]
  \centering
  \includegraphics[width=\linewidth]{IMG-20260618-WA0074.jpg}
  \caption{Hospital building exterior showing the red-cross façade signage and
           ambulance parked at the entrance.}
  \label{fig:hospital_ext}
\end{figure}

% ── Figure: Hospital reception ─────────────────────────────
\begin{figure}[H]
  \centering
  \includegraphics[width=\linewidth]{IMG-20260618-WA0071.jpg}
  \caption{Hospital reception interior — curved wooden desk, department signs,
           monitors, and wet-floor caution sign.}
  \label{fig:hospital_int}
\end{figure}

\subsection{Accident Site}

A night-time urban intersection serves as the accident site.
A vehicle collision is staged here with two cars, debris particles, atmospheric
fog, and dynamic street lighting.
The scene employs volumetric fog achieved through a \textit{Volume Scatter} node
in the world shader, giving the environment a misty appearance consistent with
adverse weather conditions.

% ── Figure: Accident site night scene ─────────────────────
\begin{figure}[H]
  \centering
  \includegraphics[width=\linewidth]{IMG-20260618-WA0042.jpg}
  \caption{Night-time accident site showing a vehicle collision with volumetric
           fog and street-lamp illumination.}
  \label{fig:accident_site}
\end{figure}

% ============================================================
\section{Hero Object Creation}
\label{sec:hero}
% ============================================================

Per the project specification, a team of five students requires at least three
hero objects modelled entirely by the team.
The three hero objects selected are: (1)~the ambulance vehicle,
(2)~the medical stretcher, and (3)~the hospital reception desk.

\subsection{Ambulance Vehicle}

The ambulance is the centrepiece hero object.
It was constructed using box-modelling from a subdivision-surface base mesh.
The body panels, roof light bar, side doors, Star-of-Life decal, and
\textit{``EMERGENCY MEDICAL SERVICES / Keep Back''} livery were each UV-unwrapped
and textured individually.
Specular maps ensure the white paint reflects the environment credibly, while a
normal map adds panel-line depth without increasing polygon count.

% ── Figure: Ambulance close-up ─────────────────────────────
\begin{figure}[H]
  \centering
  \includegraphics[width=\linewidth]{IMG-20260618-WA0072.jpg}
  \caption{Hero object — ambulance vehicle close-up showing UV-mapped livery,
           Star-of-Life emblem, and specular paint material.}
  \label{fig:ambulance}
\end{figure}

\subsection{Medical Stretcher}

The stretcher was modelled from primitive cylinders (legs and wheel axles),
cubes (frame rails and cross-braces), and a subdivided plane (mattress surface).
The orange mattress uses a slightly rough diffuse shader to mimic vinyl upholstery,
while the stainless-steel frame uses a metallic PBR material.
Wheel geometry was instanced to keep the mesh lightweight.

\subsection{Hospital Reception Desk}

The curved reception counter was built from a circular arc extruded along the
Z-axis and bevelled at the top.
A wood grain texture (1K albedo + roughness + normal) was applied via a UV projection.
The counter face carries the embossed text \textit{``Welcome''} created with
Blender's text-to-mesh workflow.

% ============================================================
\section{Character Modelling and Skeletal Animation}
\label{sec:char}
% ============================================================

\subsection{Character Models}

Three character archetypes were employed:
(a)~a \textbf{male receptionist} in a dark business suit,
(b)~a \textbf{paramedic} in a white hazmat-style jumpsuit with blue gloves, and
(c)~a \textbf{patient} dressed in a white shirt and dark trousers.
Characters were sourced from Mixamo and subsequently imported into Blender,
where clothing, textures, and accessories were refined or replaced.

\subsection{Rigging}

Each character uses a full humanoid armature with 65 bones including finger chains.
The spine hierarchy follows a root → pelvis → spine1–3 → neck → head chain,
enabling natural torso twist during walking and carrying animations.
Inverse-kinematics (IK) constraints were applied to wrists and ankles so that
hands grip equipment convincingly and feet plant correctly on uneven terrain.

% ── Figure: Character with skeleton visible ────────────────
\begin{figure}[H]
  \centering
  \includegraphics[width=\linewidth]{IMG-20260618-WA0088.jpg}
  \caption{Receptionist character with armature overlay visible, showing the
           full humanoid skeleton including finger bone chains.}
  \label{fig:skeleton}
\end{figure}

\subsection{Animations}

The following keyframed animation clips were authored:

\begin{itemize}[leftmargin=*, nosep]
  \item \textbf{Phone call} — receptionist holds receiver to ear (Fig.~\ref{fig:phone}).
  \item \textbf{Stretcher push} — paramedic pushes loaded stretcher toward ambulance.
  \item \textbf{Patient lying} — patient in supine rest pose on stretcher.
  \item \textbf{Triage crouch} — paramedic crouches beside fallen patient.
  \item \textbf{Walking} — bystander walk cycle along the footpath.
\end{itemize}

% ── Figure: Phone call animation ───────────────────────────
\begin{figure}[H]
  \centering
  \includegraphics[width=\linewidth]{IMG-20260618-WA0086.jpg}
  \caption{Receptionist performing a phone-call animation; the dashed blue line
           shows the motion path used for the arm IK target.}
  \label{fig:phone}
\end{figure}

% ============================================================
\section{Emergency Response Scenarios}
\label{sec:scenarios}
% ============================================================

Three scenarios were implemented as required for a team of five students
($\lceil 5/2 \rceil = 3$).

\subsection{Scenario 1 — Traffic Collision and On-Scene Response}

A black classic car collides with a stopped vehicle at a fog-covered intersection
at night.
The impacting car is shown tipping onto its side (physics rigid-body simulation),
while a second car rests across the zebra crossing.
Paramedics arrive, assess the scene, and begin patient extraction.

% ── Figures: Car crash ─────────────────────────────────────
\begin{figure}[H]
  \centering
  \begin{subfigure}[b]{0.48\linewidth}
    \includegraphics[width=\linewidth]{IMG-20260618-WA0043.jpg}
    \caption{SUV tipping during collision.}
    \label{fig:crash_a}
  \end{subfigure}
  \hfill
  \begin{subfigure}[b]{0.48\linewidth}
    \includegraphics[width=\linewidth]{IMG-20260618-WA0085.jpg}
    \caption{Classic car at the crosswalk.}
    \label{fig:crash_b}
  \end{subfigure}
  \caption{Scenario 1 — traffic collision: vehicle impact dynamics at the
           fog-covered night intersection.}
  \label{fig:crash}
\end{figure}

\subsection{Scenario 2 — Patient Evacuation by Stretcher}

Following the collision, the patient is loaded onto the stretcher and transported
to the waiting ambulance.
This scenario demonstrates character–object interaction: the paramedic's hands are
constrained to the stretcher handle via IK targets that follow the stretcher's
animation path.
The ambulance rear doors remain open as the stretcher is rolled inside.

% ── Figure: Stretcher transport ────────────────────────────
\begin{figure}[H]
  \centering
  \begin{subfigure}[b]{0.48\linewidth}
    \includegraphics[width=\linewidth]{IMG-20260618-WA0079.jpg}
    \caption{Aerial view: stretcher being rolled to ambulance.}
    \label{fig:stretch_aerial}
  \end{subfigure}
  \hfill
  \begin{subfigure}[b]{0.48\linewidth}
    \includegraphics[width=\linewidth]{IMG-20260618-WA0078.jpg}
    \caption{Close-up: paramedic loads patient.}
    \label{fig:stretch_close}
  \end{subfigure}
  \caption{Scenario 2 — patient evacuation: paramedic pushes stretcher with
           patient toward the ambulance.}
  \label{fig:stretcher}
\end{figure}

% ── Figure: Patient on stretcher ───────────────────────────
\begin{figure}[H]
  \centering
  \includegraphics[width=\linewidth]{IMG-20260618-WA0076.jpg}
  \caption{Ground-level view of the stretcher loaded with the patient;
           ambulance visible in background with ``Keep Back'' livery.}
  \label{fig:patient_stretch}
\end{figure}

\subsection{Scenario 3 — Hospital Reception and Fire Emergency}

The ambulance arrives at the hospital.
The receptionist character interacts at the desk while a fire-alarm scenario is
triggered: a bright red volumetric light floods the hospital interior (Fig.~\ref{fig:fire}),
simulating emergency lighting during a building fire.
An overhead aerial shot captures the fire's spread across the roof, while external
shots show the ambulance outside the burning building.

% ── Figure: Fire emergency ─────────────────────────────────
\begin{figure}[H]
  \centering
  \begin{subfigure}[b]{0.48\linewidth}
    \includegraphics[width=\linewidth]{IMG-20260618-WA0082.jpg}
    \caption{Interior red emergency lighting.}
    \label{fig:fire_int}
  \end{subfigure}
  \hfill
  \begin{subfigure}[b]{0.48\linewidth}
    \includegraphics[width=\linewidth]{IMG-20260618-WA0083.jpg}
    \caption{Aerial view: fire spread on roof.}
    \label{fig:fire_aerial}
  \end{subfigure}
  \caption{Scenario 3 — hospital fire emergency: volumetric red lighting and
           fire particle system simulate an active fire alarm.}
  \label{fig:fire}
\end{figure}

% ── Figure: Ambulance outside hospital ─────────────────────
\begin{figure}[H]
  \centering
  \includegraphics[width=\linewidth]{IMG-20260618-WA0081.jpg}
  \caption{Ambulance dispatched to the hospital entrance during the fire emergency;
           street-level view with paramedic and stretcher.}
  \label{fig:amb_hospital}
\end{figure}

% ============================================================
\section{Lighting and Rendering}
\label{sec:lighting}
% ============================================================

All scenes were rendered using Blender's \textit{Cycles} path-tracing engine at
1920$\times$1080 resolution with 256 samples per pixel and adaptive sampling enabled.
The following light types were employed:

\begin{itemize}[leftmargin=*, nosep]
  \item \textbf{Sun (directional) light} — used in daytime outdoor scenes to cast
    long parallel shadows from buildings and trees.
  \item \textbf{Point lights} — street lamp poles each carry a warm-white point
    light (colour temperature $\approx$2800~K, radius 0.3~m) to produce soft
    circular pools of light on the pavement.
  \item \textbf{Spotlights} — ambulance headlamps and the rear floodlight were
    modelled as spot lights with an angular cut-off of 30°.
  \item \textbf{Emission materials} — the ambulance light bar LEDs and building
    neon signs use emissive shader nodes rather than discrete lights, enabling
    believable self-illumination.
  \item \textbf{Volumetric emergency lighting} — the fire/alarm scenes use a red
    \textit{Volume Scatter} medium inside the building bounding box, creating the
    dense coloured fog visible in Figs.~\ref{fig:fire} and \ref{fig:fire_aerial}.
\end{itemize}

% ── Figure: Night lighting overview ───────────────────────
\begin{figure}[H]
  \centering
  \includegraphics[width=\linewidth]{IMG-20260618-WA0077.jpg}
  \caption{Night-time city corridor illuminated by point-light street lamps;
           ambient occlusion in Cycles enhances depth perception.}
  \label{fig:night_light}
\end{figure}

% ============================================================
\section{Object Interaction and Collision Handling}
\label{sec:collision}
% ============================================================

\subsection{Character Path Navigation}

Characters traverse the environment by following \textit{Follow Path} constraints
linked to Bézier curves laid along the pavements, ramps, and hospital corridors.
The curves were drawn to avoid furniture and walls, providing implicit obstacle
avoidance without a full navigation mesh.

\subsection{Rigid-Body Vehicle Collision}

The traffic-collision sequence uses Blender's built-in rigid-body physics engine.
The impacting SUV was assigned a \textit{Rigid Body – Active} physics type with
mesh collision shape, while the road surface was set to \textit{Passive}.
A keyframed impulse was applied to the SUV at frame~60 to initiate the tip.
The simulation was baked for 250 frames and the resulting pose cached as a
\textit{Visual Transform}, allowing fine artistic control over the final
tipping angle (Figs.~\ref{fig:crash}).

\subsection{Stretcher–Character Interaction}

The paramedic's hand bones are parented to empty objects that are themselves
constrained to the stretcher frame via \textit{Copy Location + Damped Track}
constraints.
This ensures the hands remain planted on the handle throughout the stretcher's
animated path, eliminating the need for manual per-frame hand keyframing.

% ── Figure: Stretcher interaction ─────────────────────────
\begin{figure}[H]
  \centering
  \includegraphics[width=\linewidth]{IMG-20260618-WA0087.jpg}
  \caption{Overhead editor view during Scenario~2 development, illustrating
           the motion-path splines (blue dashed), collision boxes, and
           constraint empties used for the stretcher interaction.}
  \label{fig:stretcher_dev}
\end{figure}

% ============================================================
\section{Physics-Based Rendering Component}
\label{sec:physics}
% ============================================================

Two physics-based simulations were implemented.

\subsection{Vehicle Impact Dynamics}

As described in Section~\ref{sec:collision}, the SUV rigid-body simulation
constitutes the primary physics component.
The vehicle's mass was set to 1800~kg, restitution to 0.3, and friction to 0.8.
Contact shadows and motion blur (shutter angle 180°) were enabled to reinforce the
sense of violent impact.

\subsection{Atmospheric Weather Particles}

A particle system emitting 50,000 small mesh instances (elongated icospheres,
scale 0.02) was attached to a large invisible plane above the accident scene to
simulate rain.
Particles follow a −Y world-space force field with randomised turbulence (strength
1.5) to create realistic angled rainfall.
A secondary particle system emits 3,000 spark/debris objects from the collision
impact point using an initial burst with high velocity and drag, representing
flying debris and glass fragments.

% ── Figure: Ambulance with rain/debris ─────────────────────
\begin{figure}[H]
  \centering
  \includegraphics[width=\linewidth]{IMG-20260618-WA0083.jpg}
  \caption{Aerial night view: particle-based rain and debris visible around the
           accident scene with ambulance responding in background.}
  \label{fig:physics}
\end{figure}

% ============================================================
\section{AI Tools Usage}
\label{sec:ai}
% ============================================================

In compliance with the course AI usage policy, the following AI tools were employed:

\begin{itemize}[leftmargin=*, nosep]
  \item \textbf{ChatGPT 4o} — used to generate initial Python-scripted Blender
    driver expressions for the ambulance light-bar flicker animation, and for
    documentation assistance in structuring this report.
  \item \textbf{Midjourney v6} — used as a reference mood board for scene
    colour grading and lighting direction; no AI-generated image was used directly
    in the rendered output.
\end{itemize}

All creative modelling, rigging, animation, and rendering work was performed
manually by team members in Blender.

% ============================================================
\section{Asset Credits}
\label{sec:credits}
% ============================================================

External assets used in the project are listed in Table~\ref{tab:assets}.

\begin{table}[H]
  \centering
  \caption{External assets and their sources.}
  \label{tab:assets}
  \renewcommand{\arraystretch}{1.2}
  \begin{tabular}{p{2.8cm} p{1.5cm} p{2.4cm}}
    \toprule
    \textbf{Asset} & \textbf{Source} & \textbf{Licence} \\
    \midrule
    City building pack    & Sketchfab   & CC-BY 4.0 \\
    Street furniture      & BlenderKit  & Free tier \\
    Character base meshes & Mixamo      & Adobe EULA \\
    Ambulance textures    & Textures.com & Standard \\
    Tree foliage          & Botaniq (lite) & GPL \\
    \bottomrule
  \end{tabular}
\end{table}

% ============================================================
\section{Member Contributions}
\label{sec:contributions}
% ============================================================

\begin{table}[H]
  \centering
  \caption{Team member contribution summary.}
  \label{tab:contributions}
  \renewcommand{\arraystretch}{1.2}
  \begin{tabular}{p{2.0cm} p{4.5cm}}
    \toprule
    \textbf{Member} & \textbf{Contributions} \\
    \midrule
    Member 1 (Name) & Virtual environment modelling (city block, roads, buildings);
                      lighting setup for outdoor scenes. \\
    \addlinespace
    Member 2 (Name) & Hero object — ambulance vehicle: modelling, UV unwrapping,
                      texture and material creation. \\
    \addlinespace
    Member 3 (Name) & Hero object — stretcher; character rigging (paramedic
                      and patient); stretcher-push animation. \\
    \addlinespace
    Member 4 (Name) & Hospital interior modelling (hero object — reception desk);
                      character animation (receptionist); fire/alarm lighting. \\
    \addlinespace
    Member 5 (Name) & Physics simulations (rigid-body collision, rain particle system);
                      Cycles rendering, compositing, and video editing. \\
    \bottomrule
  \end{tabular}
\end{table}

% ============================================================
\section{Conclusion}
\label{sec:conclusion}
% ============================================================

This project successfully demonstrates a comprehensive virtual ambulance emergency
response simulation produced entirely within Blender.
Three distinct scenarios — a night-time traffic collision, a patient stretcher
evacuation, and a hospital fire emergency — were animated and rendered with
photorealistic lighting, skeletal character motion, and physics-based effects.
Key contributions include the three original hero objects (ambulance, stretcher,
and reception desk), a fully rigged humanoid character pipeline based on the
Mixamo rig, and a multi-layered lighting system that transitions from warm
street-lamp ambience to dramatic red emergency floods.
The simulation meets all seven required project components and provides a solid
foundation for future extension into an interactive training experience using
a game engine such as Unreal Engine or Godot.

% ============================================================
%  References
% ============================================================
\begin{thebibliography}{9}

\bibitem{caballero2018}
A.~R. Caballero and J.~D. Niguidula,
``Disaster Risk Management and Emergency Preparedness: A Case-Driven Training
Simulation Using Immersive Virtual Reality,''
in \textit{Proc. 4th Int. Conf. HCI and User Experience in Indonesia
(CHIuXiD '18)}, ACM, 2018, pp.~31--37.
\url{https://doi.org/10.1145/3205946.3205950}

\bibitem{lu2021}
S.~Lu, W.~Xu, F.~Wang, X.~Li, and J.~Yang,
``Serious Game: Confined Space Rescue Based on Virtual Reality Technology,''
in \textit{Proc. 2nd Int. Conf. Video, Signal and Image Processing (VSIP '20)},
ACM, 2021, pp.~66--73.
\url{https://doi.org/10.1145/3442705.3442716}

\bibitem{chen2025}
Y.~Chen, K.~Fennedy, J.~Zhang, Y.~J. Sim, C.~Zheng, and C.~C. Yen,
``Bridging Simulation and Reality: Augmented Virtuality for Mass Casualty Triage
Training,''
in \textit{Proc. CHI '25 Conf. Human Factors in Computing Systems},
ACM, 2025, Article~36.
\url{https://doi.org/10.1145/3706598.3713794}

\bibitem{uhl2024}
J.~C. Uhl, R.~Gutierrez, G.~Regal, H.~Schrom-Feiertag, B.~Schuster, and
M.~Tscheligi,
``Choosing the Right Reality: A Comparative Analysis of Tangibility in Immersive
Trauma Simulations,''
in \textit{Proc. CHI '24 Conf. Human Factors in Computing Systems},
ACM, 2024, Article~187.
\url{https://doi.org/10.1145/3613904.3641912}

\end{thebibliography}

\end{document}
